\subsection{线性群中的指数映射}

命题 17. 设 \(\Delta\) 是由 \(z\in \mathbf{C}\) 满足 \(-\pi < \mathcal{I}(z) < \pi\) 所成的集合，而 \(\Delta'\) 是由 \(z\in \mathbf{C}\) 不是实数 \(\leqslant 0\) 所成的集合。设 \(E\) 是 \(\mathbf{C}\) 上的完备可赋范空间，且 \(A\)（分别地，\(A'\)）是由 \(x\in \mathcal{L}(E)\) 且其谱 \(\mathrm{Sp}(x)\) 包含于 \(\Delta\)（分别地，\(\Delta'\)）的元素构成的集合。则 \(A\)（分别地，\(A'\)）是 \(\mathcal{L}(E)\)（分别地，\(\mathbf{GL}(E)\)）的开子集，且映射 \(\exp: A\to A'\) 和 \(\log: A'\to A\)（谱理论，第 I 章，§4，第 9 号）是解析流形的互逆同构。

这由谱理论，第 I 章，§ 4，命题 10 和第 9 号得到。

定理 6. 设 E 是实或复 Hilbert 空间，U 是 E 的酉群。

\begin{enumerate}
\def\labelenumi{(\roman{enumi})}
\item
  \(\mathbf{H}\) 是 \(\mathcal{L}(\mathbf{E})\) 中厄米元素的集合；按实赋范空间结构，它是 \(\mathcal{L}(\mathbf{E})\) 的闭向量子空间，并具有拓扑补空间。
\item
  \(\mathbf{H}'\) 中满足 \(\geqslant 0\) 的元素的集合 \(\mathbf{GL}(\mathbf{E})\) 是 \(\mathbf{GL}(\mathbf{E})\) 的实解析子流形。
\item
  映射 exp 在 H 上的限制是从 H 到 H' 的实解析流形同构。
\item
  从 \((h,u)\mapsto (\exp h)u\) 到 \(\mathbf{H}\times \mathbf{U}\) 的映射 \(\mathbf{GL}(\mathbf{E})\) 是实解析流形同构。
\end{enumerate}

回忆，若 \(x \in \mathcal{L}(\mathrm{E})\)，则 \(x^*\)\(x\) 表示 x 的伴随。令 \(\mathrm{H}_1\) 为满足 \(x \in \mathcal{L}(\mathrm{E})\) 的 \(x^* = -x\) 的集合。公式 \(x = \frac{1}{2}(x + x^*) + \frac{1}{2}(x - x^*)\) 说明，按其实赋范空间结构，\(\mathcal{L}(\mathrm{E})\) 是 \(\mathrm{H}\) 与 \(\mathrm{H}_1\) 的拓扑直和，从而得到 (i)。

假设 \(\mathbf{K} = \mathbf{C}\)。在命题 17 的记号下，\(\mathrm{H}'\) 是满足 \(h \in \mathrm{H} \cap \Delta'\) 的 \(\mathrm{Sp} h \subset \mathbf{R}_{+}^{*}\) 的集合。由于 \(\exp(\mathbf{R}) = \mathbf{R}_{+}^{*}\)，(ii) 和 (iii) 由命题 17 以及谱理论，第 I 章，§4，命题 8 和 §6，第 5 号得到。从 \((h, u) \mapsto y = (\exp h)u\) 到 \(\mathrm{H} \times \mathrm{U}\) 的映射 \(\mathbf{GL}(\mathrm{E})\) 是双射，这是谱理论，第 I 章，§6，命题 15 的结论。由上可知它是实解析的。映射 \(y \mapsto h = \frac{1}{2} \log(yy^*)\) 是实解析的，因而映射 \(y \mapsto u = (\exp h)^{-1}y\) 也是实解析的。故得 (iv)。

假设 \(\mathbf{K} = \mathbf{R}\)。令 \(\hat{\mathbf{E}}\) 为 E 的复化 Hilbert 空间，令 J 为将 \(\xi + i\eta \mapsto \xi - i\eta\) 的映射（其中 \(\xi, \eta\) 属于 E），它从 \(\hat{\mathbf{E}}\) 到 \(\hat{\mathbf{E}}\)。令 \(\tilde{\mathbf{H}}, \tilde{\mathbf{H}}', \tilde{\mathbf{U}}\) 表示对 \(\hat{\mathbf{E}}\) 定义的集合，正如 H、H'、U 对 E 定义的那样。于是 H（分别地，H'、U）可与满足 \(x \in \tilde{\mathbf{H}}\) 的 \(\tilde{\mathbf{H}}'\)（分别地，\(\tilde{\mathbf{U}}\)、\(\mathrm{JxJ}^{-1} = x\)）的集合相识别。性质 (ii)、(iii) 和 (iv) 随即由 (i) 及复情形中的相应性质容易得到。

命题 18. 设 E 是 C 上的完备可赋范空间，\(v \in \mathcal{L}(E)\)，且 g = exp v。假定 \(\operatorname{Sp}(v)\) 不包含形如 \(2i\pi n\) 的点，其中 \(n \in Z - \{0\}\)。则对所有 \(x \in E\)，条件 vx = 0 与 gx = x 等价。

这由第 4 号的引理 2 应用于函数 \(z \mapsto e^z - 1\) 得到。

推论 1. 设 E 是 C 上的完备可赋范空间，F 是从 \(\mathbf{E}^n\) 到 E 的连续 n-线性映射空间。对所有 \(v \in \mathcal{L}(\mathbf{E})\)，令 \(\sigma(v)\) 为由下式定义的 \(\mathcal{L}(\mathbf{F})\) 中的元素：

\[
(\sigma (v) f) (x _ {1}, \dots , x _ {n}) = v (f (x _ {1}, \dots , x _ {n})) - \sum_ {i = 1} ^ {n} f (x _ {1}, \dots , v x _ {i}, \dots , x _ {n}).
\]

对所有 \(g \in \mathbf{GL}(\mathbf{E})\) ，令 \(\rho(g)\) 为由下式定义的 \(\mathbf{GL}(\mathbf{F})\) 中的元素：

\[
(\rho (g) f) (x _ {1}, \dots , x _ {n}) = g (f (g ^ {- 1} x _ {1}, \dots , g ^ {- 1} x _ {n})).
\]

设 \(u \in \mathcal{L}(\mathrm{E})\) 使得每个 \(z \in \operatorname{Sp} u\) 都满足 \(|\mathcal{I}(z)| < \frac{2\pi}{n + 1}\)。则对所有 \(f \in \mathrm{F}\)，条件 \(\sigma(u)f = 0\) 与 \(\rho(\exp u)f = f\) 等价。

\(L(\rho) = \sigma (\S 3, \text{no. } 11, \text{命题 41 的推论 1}) \text{从而}\)

\[
\rho (\exp u) = \exp \sigma (u)
\]

（第 4 号，命题 10 的推论 3）。由命题 18 可知，只需证明 \(\operatorname{Sp} \sigma(u)\) 不与 \(2i\pi(\mathbf{Z} - \{0\})\) 相交。但这由下述引理得到：

引理 6. 若 \(v \in \mathcal{L}(\mathrm{E})\)，则 \(\operatorname{Sp} \sigma(v) \subset \operatorname{Sp} v + \operatorname{Sp} v + \cdots + \operatorname{Sp} v\)，其中和式含有 \(n + 1\) 项。

对所有 \(v_{0}, v_{1}, \ldots, v_{n}\)，我们通过下式定义 \(\mathcal{L}(\mathrm{F})\) 中的元素 \(f \in \mathrm{F}\)：

\[
\begin{array}{l} (v _ {0} f) (x _ {1}, \dots , x _ {n}) = v (f (x _ {1}, \dots , x _ {n})) \\ (v _ {i} f) (x _ {1}, \dots , x _ {n}) = - f (x _ {1}, \dots , v x _ {i}, \dots , x _ {n}) \quad \text { for } 1 \leqslant i \leqslant n. \end{array}
\]

于是 \(\sigma(v) = \sum_{i=0}^{n} v_i\)，且各 \(v_i\) 两两可交换。令 A 为由 \(\mathcal{L}(F)\) 生成的 \(v_i\) 的闭子代数；它是交换的（

谱理论，第 I 章，§ 1，第 4 号），且 \(\mathrm{Sp}_{\mathcal{L}(\mathbb{F})}v' = \mathrm{Sp}_A v' \subset \sum_{i=0}^{n} \mathrm{Sp} v_i\)（谱理论，第 I 章，§ 3，命题 3 (ii)）。现在，若 \(\lambda \in \mathbf{C}\) 使得 \(v - \lambda\) 可逆，则显然 \(v_i - \lambda_i\) 可逆，从而对所有 \(\mathrm{Sp} v_i \subset \mathrm{Sp} v\)，\(i\)。

推论 2. 设 \(\mathbf{E}\) 是 \(\mathbf{C}\) 上的完备可赋范代数，且 \(w \in \mathcal{L}(\mathbf{E})\)。假定每个 \(z \in \operatorname{Sp} w\) 都满足 \(|\mathcal{I}(z)| < \frac{2\pi}{3}\)。下列条件等价：

\begin{enumerate}
\def\labelenumi{(\roman{enumi})}
\item
  \(w\) 是 \(\mathbf{E}\) 的导子；
\item
  exp w 是 E 的自同构。
\end{enumerate}

这由取 \(n = 2\) 且令 \(f\) 为 E 的乘法的推论 1 得到。

命题 19. 设 E 是 C 上的完备可赋范空间，\(v \in \mathcal{L}(E)\)，且 g = exp v。假定每个 \(z \in Sp\) v 都满足 \(-\pi < \mathcal{I}(z) < \pi\)。则对于 E 的每个闭向量子空间 \(E'\)，条件 \(v(E') \subset E'\) 与 \(g(E') = E'\) 等价。

条件 \(v(\mathrm{E}') \subset \mathrm{E}'\) 蕴含 \(g(\mathrm{E}') \subset \mathrm{E}'\) 以及 \(g^{-1}(\mathrm{E}') \subset \mathrm{E}'\)，因而 \(g(\mathrm{E}') = \mathrm{E}'\)。假设 \(g(\mathrm{E}') = \mathrm{E}'\)。我们采用命题 17 的记号 \(\Delta, \Delta'\)。由于 Sp \(v\) 是 \(\Delta\) 的紧子集，存在紧矩形 \(\mathbf{Q} = [\underline{a}, b] \times [\underline{a}', b']\)，使得 \(\operatorname{Sp} v \subset \mathbf{Q} \subset \Delta\)。集合 \(\Delta - \mathbf{Q}\) 是连通的。因此 \(\operatorname{Sp} g \subset \exp \mathbf{Q} \subset \Delta'\)，集合 \(\exp \mathbf{Q}\) 是紧的，而集合 \(\Delta' - \exp \mathbf{Q}\) 是连通的。后者的闭包包含 \(] - \infty, [0]\)，从而

\[
\left. \left(\Delta^ {\prime} - \exp Q\right) \cup ] - \infty , 0 ] = C - \exp Q \right.
\]

是连通的。于是  是多项式凸的（谱理论，第 I 章，§3，命题 9 的推论 2），从而定义在 \(\Delta'\) 上的函数 log 是 \(\mathcal{O}(\exp Q)\) 中多项式函数的极限（谱理论，第 I 章，§4，命题 3）。因此 \(v = \log g\) 是 \(\mathcal{L}(E)\) 中形如 P(g) 的元素的极限，其中 P 是多项式（谱理论，第 I 章，§4，定理 3）。由于 P(g)(E') \(\subset\) E'，可得 \(v(E') \subset E'\)。

推论。设 E 是 C 上的完备赋范空间，\(v \in \mathcal{L}(E)\)，且 \(g = \exp v\)。

假定每个 \(z \in \operatorname{Sp} v\) 都满足 \(-\frac{\pi}{2} < \mathcal{I}(z) < \frac{\pi}{2}\)。那么，对于每个闭向量

子空间 M 的条件 \(\mathcal{L}(\mathrm{E})\)\(g\mathrm{M}g^{-1} = \mathrm{M}\) 与 \([v, \mathrm{M}] \subset \mathrm{M}\) 等价。

令 \(\mathrm{F} = \mathcal{L}(\mathrm{E})\)，令 \(g'\) 为从 \(f \mapsto gfg^{-1}\) 到 \(\mathrm{F}\) 的映射 \(\mathrm{F}\)，令 \(v'\) 为从 \(f \mapsto [v, f]\) 到 \(\mathrm{F}\) 的映射 \(\mathrm{F}\)。则 \(g' = \exp v'\)（第 4 号，命题 10 的推论 3 以及 §3，第 11 号，命题 41 的推论 1）。引理 6 说明，对所有 \(-\pi < \mathcal{I}(z) < \pi\)，都有 \(z \in Sp v'\)。于是只需应用命题 19。

